\section{Noise Robustness}
\label{sec:noise}

This section is intentionally short: explicit noise-robustness training
is not part of the v3 checkpoint described in this report. We document
the planned evaluation protocol so that the next release can be
measured against this baseline.

\subsection{Planned target benchmark}

The chosen reference benchmark for the next iteration is
\textbf{Whisper-RIR-Mega test}, a far-field / reverberant test split
that mixes synthetic room impulse responses (RIRs) with real-world
noise. The dataset is already archived internally
% TODO: confirm the canonical citation / source URL for
% Whisper-RIR-Mega test and add the corresponding bib entry.
under
\texttt{/ai\_sds\_wuzz/MULTILINGUAL\_DATA/noise/rirmega/data/manifests}.

\subsection{Reference baseline}

The reference baseline is Qwen3-ASR-1.7B \cite{qwen2025qwen3asr},
chosen because (i)~it is a strong open-source small model with
documented robustness on noisy / far-field audio
\cite{qwen2025qwen3asr}, and (ii)~it occupies the same family of
LLM-ASR designs as AmphionASR. The bar for the next release is to
match-or-beat Qwen3-ASR-1.7B's WER on Whisper-RIR-Mega test, in
particular at low SNRs.

\subsection{Planned recipe}

The plan is to add a noise-augmented mixture (clean speech convolved
with internal-and-public RIR sets, with stationary and non-stationary
noise mixed at SNRs covering a range of low-to-medium values) to
Stage~2 / Stage~3 training, while keeping the rest of the recipe
unchanged. Results will appear in a follow-up revision of this
technical report.
% TODO: replace this section with concrete numbers as soon as the
% noise-augmented checkpoint exists.
